Graph-structured data have become a fundamental component of modern information systems, including social networks, recommendation engines, knowledge graphs, transportation networks, and communication infrastructures. As the size of these graphs continues to grow, graph databases increasingly rely on distributed architectures in which graph data are partitioned across multiple storage nodes. While distribution enables horizontal scalability, it also introduces the challenge of efficiently assigning graph vertices to partitions.

The quality of graph partitioning has a direct impact on the performance of distributed graph databases. During graph traversals and shortest-path queries, communication between storage nodes is required whenever adjacent vertices reside in different partitions. Such inter-partition communication incurs significantly higher latency than local memory access and can become the dominant factor limiting system performance. Consequently, an effective partitioning strategy should simultaneously minimize the graph cut while maintaining a balanced distribution of data across storage nodes.

Numerous graph partitioning techniques have been proposed to address this problem. Some graph databases, such as ArangoDB \cite{arangodb_github}, prioritize balanced data distribution using hash-based partitioning rather than optimizing partition quality. Traditional partitioning frameworks, such as METIS \cite{MetisOG}, produce high-quality partitions by exploiting the structural properties of the graph. However, these methods are primarily designed for static graphs and often require global repartitioning, making them computationally expensive for large-scale or dynamically evolving datasets. In contrast, streaming partitioning algorithms, such as Fennel \cite{Fennel}, provide scalable and efficient vertex placement as the graph is constructed, balancing partition sizes while reducing the expected graph cut. Despite their computational efficiency, streaming methods are inherently greedy and may produce suboptimal partitions that cannot be improved without additional refinement. Furthermore, because they make decisions without knowledge of the complete graph structure, their partition quality may deteriorate when vertices are processed in an unfavorable order.

This work investigates the potential of a hybrid graph partitioning approach that combines the scalability of the Fennel streaming algorithm with the local optimization capabilities of the Kernighan--Lin (KL) algorithm. In the proposed method, Fennel is employed to construct an initial balanced partitioning of the graph, after which KL refinement is applied to reduce the number of inter-partition edges while preserving partition balance. The proposed combination aims to achieve a favorable trade-off between partitioning efficiency and partition quality, making it suitable for large distributed graph databases where both scalability and communication cost are critical considerations.

The hypothesis of this work is that Fennel provides an initial partitioning of sufficient quality such that a small number of KL refinement iterations can achieve higher partition quality while substantially reducing computational cost compared to the complete METIS partitioning process.

